% Bevi -- from a web-only Expo app to a native iOS application: a migration study.
% Build: tectonic docs/bevi-ios-migration.tex   (or pdflatex twice)
% Written 2026-11-14 against be/custom-bev-icons @ 31381bc (contract 3.0.0, Expo SDK 57.0.31,
% React Native 0.86.3). Every path, count and configuration quoted was measured in the tree on
% that date. Nothing in the repository was modified to produce this document.

\documentclass[11pt]{article}

\usepackage[a4paper,margin=1.05in,bottom=1.1in]{geometry}
\usepackage{libertinus}
\usepackage{inconsolata}
\usepackage{microtype}
\usepackage{booktabs}
\usepackage{array}
\usepackage{enumitem}
\usepackage{xcolor}
\usepackage{titlesec}
\usepackage{parskip}
\usepackage[most]{tcolorbox}
\usepackage{fancyhdr}
\usepackage{needspace}
\usepackage{etoolbox}
\usepackage[colorlinks=true,linkcolor=ink,urlcolor=ink]{hyperref}

% ---------------------------------------------------------------- palette ---
\definecolor{ink}{HTML}{1F3A5F}      % headings and rules
\definecolor{accent}{HTML}{B4540A}   % the one warm accent, for warnings
\definecolor{wash}{HTML}{F4F6F9}     % callout background
\definecolor{edge}{HTML}{C9D4E2}     % callout border
\definecolor{mute}{HTML}{5A6672}     % captions, metadata

% ------------------------------------------------------------- typography ---
\linespread{1.08}
\setlength{\parskip}{0.55em}
\setlist{itemsep=0.45em,parsep=0pt,topsep=0.35em,leftmargin=1.35em}
\newcommand{\code}[1]{\texttt{\small #1}}
\newcommand{\bk}{\discretionary{}{}{}}
\tolerance=1800
\emergencystretch=3em
\hbadness=10000

\titleformat{\section}{\color{ink}\Large\bfseries}{\color{ink}\thesection}{0.7em}{}[\vspace{0.15em}{\color{edge}\titlerule[1pt]}]
\titlespacing*{\section}{0pt}{1.9em plus .3em}{0.9em}
\titleformat{\subsection}{\color{ink}\large\bfseries}{\color{ink}\thesubsection}{0.6em}{}
\titlespacing*{\subsection}{0pt}{1.5em plus .2em}{0.55em}

% Inventory entries: the file on its own line, the explanation beneath it.
\newenvironment{findings}
  {\begin{description}[style=nextline,leftmargin=1.1em,labelindent=0pt,
     itemsep=1.05em,parsep=0.3em,topsep=0.6em,font=\normalfont\color{ink}\bfseries]}
  {\end{description}}

\newtcolorbox{callout}[1][]{enhanced,colback=wash,colframe=edge,
  boxrule=0.6pt,left=1.1em,right=1.1em,top=0.9em,bottom=0.9em,arc=2pt,#1}

\newtcolorbox{warnbox}[1][]{enhanced,colback=white,colframe=accent,
  boxrule=0.9pt,left=1.1em,right=1.1em,top=0.9em,bottom=0.9em,arc=2pt,#1}

\widowpenalty=10000
\clubpenalty=10000
\pretocmd{\section}{\needspace{6\baselineskip}}{}{}
\pretocmd{\subsection}{\needspace{4.5\baselineskip}}{}{}

\pagestyle{fancy}
\fancyhf{}
\renewcommand{\headrulewidth}{0pt}
\fancyfoot[C]{\color{mute}\small\thepage}

\begin{document}

% ------------------------------------------------------------------ head ---
\begin{center}
  {\color{ink}\Huge\bfseries Bevi on iOS\par}
  \vspace{0.7em}
  {\large From a web-only Expo app to a native iOS application\par}
  \vspace{0.45em}
  {\color{mute}\large a migration study of the current codebase\par}
  \vspace{1.6em}
  {\color{mute}\small Bevi engineering --- Arjun's and Matt's lanes \quad\textbullet\quad 14 November 2026\par}
  \vspace{0.3em}
  {\color{mute}\small measured against \code{be/custom-bev-icons} @ \code{31381bc}, contract 3.0.0\par}
\end{center}

\vspace{1.4em}

\begin{callout}
\textbf{\color{ink}The short version}

\begin{itemize}[leftmargin=1.2em,itemsep=0.5em,topsep=0.5em]
  \item Bevi is \emph{already} a React Native application. \code{app/} is Expo SDK 57.0.31 on React
  Native 0.86.3 with Expo Router, deliberately configured web-only on 2026-10-11 and shipped as a
  single-page \code{react-native-web} build on Vercel.
  \item Buying the Apple Developer Program therefore buys a build target that was never switched on.
  This is not a rewrite. The work is repaying the web-only assumptions that decision accumulated.
  \item The backend needs no change at all --- not one migration, RLS policy or Edge Function.
  \item Two items are genuine blockers rather than work items: the \code{Blob}-typed media contract,
  and the absence of user blocking under App Review guideline~1.2.
  \item Realistic effort: \textbf{four to six weeks} of app-side work before a TestFlight build we
  would be happy to show people.
\end{itemize}
\end{callout}

\vspace{0.6em}

% --------------------------------------------------------------------- 1 ---
\section{Where we stand today}

\code{app/src} holds 305 TypeScript and TSX files, 30{,}840 lines, 47 design-system components and
103 test files. All of it is written against React Native primitives --- \code{View},
\code{Pressable}, \code{StyleSheet}, Reanimated, \code{FlashList}, \code{expo-image} --- rather than
DOM elements. The web build is \code{react-native-web} translating those primitives into HTML, not a
separate web codebase.

Four earlier decisions leave us unusually well placed, by design rather than luck.

\subsection{The backend needs no change}

Auth is username and password over synthetic emails (\code{usernameToEmail},
\code{@users.bevi.app}). That means no OAuth redirect, no magic link and no callback URL to register
anywhere. \code{packages/api/src/client.ts} already sets \code{detectSessionInUrl: false}, which is
the native-correct setting. The four Edge Functions carry no origin allow-list, the 60 migrations and
every RLS policy are origin-agnostic, and only the publishable key ever reaches a client --- exactly
as today.

\subsection{The seams we need already exist}

\code{createBeviApi} takes an injected \code{AuthStorageAdapter} whose methods may return promises:
precisely the shape AsyncStorage needs, and
\code{@react-native-\bk async-storage/\bk async-storage} is already a dependency.
\code{app/.gitignore} already ignores \code{ios/} and \code{android/} as ``generated by
\code{expo prebuild}'', so there is no checked-in Xcode project to reconcile, ever.

\subsection{No typography work}

\code{src/theme/tokens.ts} is system-font only and is parity-tested against the design system, so
there are no fonts to load and the visual contract crosses platforms untouched.

\subsection{Three installed libraries are degraded today}

\code{expo-haptics} is a no-op on web. \code{expo-glass-effect} never renders real glass.
\code{expo-image-picker} falls back to an \code{<input type="file">}. All three are already installed
and already wired into screens; on a device they simply become real, at no integration cost.

% ------------------------------------------------------------------- 2 ---
\section{Is this actually worth doing?}

The \$99 is not the decision. The decision is four to six weeks of app-lane time and, permanently,
accepting a third party between \code{main} and production. Here is the honest case on both sides.

\subsection{What we actually gain}

\textbf{Push notifications --- the biggest single item.} Bevi has none today, and nothing in TASKS.md
plans any, because the web build effectively cannot have them: iOS delivers web push only to an
installed Home Screen app and does so unreliably. Bevi is a social product whose loop is follows,
feeds, comments and other people's rankings, and that loop currently has no way to reach anyone who
is not already looking at the app. Natively this is \code{expo-notifications} plus an APNs key.

\textbf{The install path.} Today the only way in is Safari, then the share sheet, then ``Add to Home
Screen'' --- a flow most people have never used and will not discover. Growth is invite-driven, so we
pay that conversion tax at exactly the moment an invite lands in someone's hand.

\textbf{Sessions that survive.} Safari evicts script-writable storage for sites that are not
installed, so \code{localStorage} sessions quietly disappear and people find themselves signed out
for no reason they can see. On a device the session lives in app storage and simply persists.

\textbf{Find Bev stops fighting the browser.} The reviewed plan in
\code{docs/find-bev-in-\bk app-plan.md} is shaped end to end by WebKit's behaviour in a Home Screen
app: one URL for the whole trip, a \code{Modal} rather than a route, and an iPhone spike (S2) that
gates the feature. Natively those constraints evaporate, and camera, motion and location become
ordinary libraries.

\textbf{The feel we already designed for.} \code{expo-glass-effect} and \code{expo-haptics} are
installed and inert. Native navigation gestures and transitions, the system share sheet and photo
picker, Reanimated on the UI thread and a genuinely recycling \code{FlashList} are all things the web
build approximates and a device simply has.

\textbf{Surfaces that are closed today:} widgets, Live Activities, Shortcuts, App Clips. None are on
the roadmap, but none are reachable from a web build either.

\subsection{What it costs}

\textbf{Time, and what that time is not spent on.} Four to six weeks of app-lane work is roughly the
Map tab plus the duplicates work. That opportunity cost, not the membership fee, is the real price.

\textbf{App Review becomes a release gate.} Today a merge to \code{main} is in production within
minutes and nobody outside the two of us is involved. Natively, every release waits on review --- a
day or two when it goes well --- and guideline 1.2 work has to land before the first submission at
all. For a two-person team that ships several times a week, this is a genuine change in tempo.

\textbf{Version skew becomes a permanent discipline, and this is the most underrated cost.} Today
every client runs \code{main} and \code{CONTRACT\_VERSION} moves in lockstep with the backend. Once
builds live in the App Store, old clients persist for weeks and the contract has to tolerate them.
The contract already anticipates this --- \code{AppConfig} carries \code{minAppVersion} and
\code{maintenance} --- but \code{createMetaApi} returns a fixed placeholder, with the comment that no
versioning or maintenance-flag infrastructure exists yet. That infrastructure has to become real, and
contract changes afterwards must either stay backward-compatible or be gated behind a minimum
version. It belongs in phase 1, not ``later''.

\textbf{The accessibility sweep is invisible work.} Eighty-three ARIA props that no test will catch
and no screenshot will show. It is exactly the kind of task that gets deferred and then quietly never
happens.

\textbf{Either a dual-target tax or a lost surface.} Keeping the web build means a
\code{Platform} branch or a file split for every native-only API, forever, and two QA passes.
Dropping it means losing shareable profile links, the browser as a demo and bug-reproduction surface,
and the Playwright smoke path.

\textbf{Android becomes a separate decision.} Today the web build covers Android for free. A native
iOS app does not, and an Android build is a second review process, a second store listing and a
second device matrix --- cheap from the same Expo codebase, but not free.

\textbf{Small, real overheads:} \$99 a year, a Mac with Xcode, EAS queue time or a paid tier, and the
8.2\,MB icon question.

\subsection{What does not change --- why the bet is bounded}

The backend is untouched: 60 migrations, every RLS policy, the four Edge Functions, the ranking
package, the design system, the catalog and all 103 test files carry over exactly as they are. The
move lives inside one lane, and the web build can stay alive the entire time. That makes this
reversible in a way a rewrite would not be --- if phase 1 goes badly, we still have the app we have
today, deploying to Vercel on every merge.

\subsection{The counter-case, stated fairly}

If the next quarter is catalog work, ranking, venues and fast iteration on the web, then App Review
is pure drag and the PWA already does nearly everything --- everything except notifications and
reliable sensors. The cheap path is to keep shipping web, run spike S2 as planned, and revisit this
when notifications actually block something. That is a legitimate choice, and it costs nothing today.

\begin{callout}
\textbf{\color{ink}Verdict}

\begin{itemize}[leftmargin=1.2em,itemsep=0.5em,topsep=0.5em]
  \item \textbf{Buy the membership now regardless.} \$99 is noise next to the time already going into
  this question, and it immediately unblocks running on a real device --- which de-risks spike S2 on
  either path.
  \item \textbf{Treat phase 0 as a two-day spike, not a commitment.} If the app comes up on a phone
  and the main flows behave, everything after it is grind with few unknowns. If it does not, we have
  spent two days.
  \item \textbf{Commit to phases 1--4 only if notifications and camera-based Find Bev are genuinely
  on the roadmap for the next two quarters.} They are the two things the web build cannot do. If they
  are not, stay on web and revisit without regret.
  \item \textbf{If we go, build the \code{minAppVersion} gate in phase 1}, before the first build
  reaches anyone outside the team.
\end{itemize}
\end{callout}

% ------------------------------------------------------------------- 3 ---
\section{The Apple-side work}

\textbf{Enrolment.} The Apple Developer Program is \$99 per year. Enrol as an \emph{organization}
only if the seller name must read as a company --- that path needs a D-U-N-S number and adds days.
Individual enrolment is immediate but publishes a personal name as the seller.

\textbf{Identifiers.} Register the bundle identifier (it is permanent and can never be reused), then
create the App Store Connect record.

\textbf{Configuration.} Add an \code{ios} block to \code{app/app.config.ts}, which today holds only a
\code{web} key and the comment ``No native builds, so no \code{ios}/\code{android} blocks, bundle ids
or EAS profiles.'' That block must carry:

\begin{itemize}
  \item \code{bundleIdentifier} and \code{buildNumber};
  \item a \code{scheme}, for deep links;
  \item a 1024\,px icon and an \code{expo-splash-screen} configuration;
  \item the \code{expo-image-picker} plugin's camera and photo-library permission strings;
  \item \code{ios.privacyManifests}, for the \code{PrivacyInfo.\bk xcprivacy} required-reason
  declarations Apple has required since 2024;
  \item \code{ITSAppUses\bk NonExemptEncryption: false} (we are HTTPS only).
\end{itemize}

\textbf{Building.} Either locally with \code{npx expo run:ios} (Xcode on a Mac) or on EAS Build via a
new \code{eas.json} with development, preview and production profiles. The two
\code{EXPO\_\bk PUBLIC\_\bk SUPABASE\_*} variables move from Vercel's project settings into those
build profiles, where they are inlined into the IPA exactly as they are inlined into the web bundle
today. Nothing secret ships either way.

\textbf{Distribution.} TestFlight: up to 100 internal testers immediately, and 10{,}000 external
testers once the first external build passes a review.

% --------------------------------------------------------------------- 3 ---
\section{Inventory: every web-only assumption we have to repay}
\label{sec:inventory}

Eleven non-test files branch on \code{Platform.OS}; eighteen touch browser globals. What follows is
the complete list, grouped by theme. Each entry names the file, what it assumes, and what native
needs instead.

\subsection{Session and storage}

\begin{findings}

\item[\code{app/src/api/authStorage.ts}]
\code{webAuthStorage()} wraps \code{localStorage}, with an in-memory fallback for private-mode
Safari. Native needs an AsyncStorage (or \code{expo-secure-store}) adapter behind the same interface
--- a small, well-isolated change, because the interface is already injected. It also needs
\code{auth.\bk startAutoRefresh} and \code{stopAutoRefresh} driven by React Native's \code{AppState}:
native has no tab-visibility signal, so without this a backgrounded app comes back with a stale
token.

\item[\code{features/add-bev/recents.ts}, \code{features/bug-report/draft.ts},
\code{lib/clearLegacy\bk Storage.ts}]
Three further readers and writers of \code{localStorage}, all of them \emph{synchronous}.
AsyncStorage is asynchronous, so these call sites change shape rather than merely swapping their
backend. Budget for the call-site churn, not just the adapter.

\end{findings}

\subsection{Media and photos}
\label{sec:media}

\begin{findings}

\item[\code{shared/src/api.ts:373} and \code{packages/api/src/media.ts}]
\code{UploadMediaInput.\bk file} is typed \code{Blob}, and that blob is handed straight to
\code{storage.upload()}. React Native's \code{Blob} is not a real byte container, and Supabase
Storage expects an \code{ArrayBuffer} or \code{FormData} from a device. Widening the contract is a
\code{CONTRACT\_VERSION} bump touching both lanes, \code{FakeApi} and the conformance suite. See the
blocker note below.

\item[\code{features/create-bev/photo.ts}]
Built on \code{URL.createObjectURL}, \code{blob:} preview URIs and a
\code{document.createElement('canvas')} re-encode fallback. The good news:
\code{expo-image-\bk manipulator} works natively and returns \code{file://} URIs, so the canvas path
becomes dead code and \code{revokeIfBlob} a no-op. The thing to watch: this file carries a privacy
guarantee --- an image that cannot be decoded is \emph{rejected} rather than uploaded untouched, so
GPS EXIF is always stripped --- and that guarantee has to be re-proved on the native path.

\end{findings}

\subsection{Assets and icons}

\begin{findings}

\item[\code{app/src/lib/bevIcons.\bk generated.ts} --- 192 SVGs, 8.2\,MB in
\code{app/public/bev-icons/}]
Every value is an absolute web path (\code{'/bev-icons/x.svg?v=hash'}), and Expo's \code{public/}
directory is a web-only feature that Metro never bundles into a native app. Two options: a generated
\code{require()} asset map, regenerating \code{scripts/icons/\bk export-app.mjs} and adding roughly
8\,MB to the IPA; or remote URLs served from Supabase Storage, which keeps the binary small at the
cost of a cold-launch cache. Either way, confirm \code{expo-image}'s SVG decoding on a real device at
56\,pt --- \code{react-native-svg} is already installed as a fallback.

\end{findings}

\subsection{Accessibility --- the largest invisible regression}

\begin{findings}

\item[83 \code{aria-*} props across 20 non-test files, most of them \code{aria-pressed}]
\code{react-native-web} renders ARIA attributes, and several of ours are documented workarounds for
it dropping \code{accessibilityState}. iOS ignores every one of them, silently: nothing fails, no
test goes red, and VoiceOver simply stops announcing toggle state. Each needs its React Native
equivalent --- \code{accessibilityState}, \code{accessibility\bk Value}, \code{accessibilityRole} ---
verified by hand under VoiceOver. This is the single largest item in the inventory and the easiest to
under-budget.

\end{findings}

\subsection{Layout, chrome and platform polish}

\begin{findings}

\item[\code{components/BevGlow.tsx}, \code{FeedPost.tsx}, \code{CommentsSheet.tsx},
\code{Sheet.tsx}]
Web-only style props: \code{filter: blur(48px)}, \code{touchAction}, \code{cursor},
\code{overflowWrap}, and a \code{SAFE\_BOTTOM} constant computed for the browser. React Native 0.76
and later, on the New Architecture, support \code{filter} and \code{boxShadow} natively, so most of
this survives --- but verify prop by prop on a device. Safe-area insets come from
\code{react-native-\bk safe-area-context} instead of CSS environment variables.

\item[\code{components/BucketButton.tsx} and \code{lib/glass.ts}]
TASKS.md records that ``no Apple developer route is needed on the web build.'' That caveat now
expires. \code{GlassView} renders real glass on iOS 26 and later, and Reduce Transparency must be
read through \code{AccessibilityInfo} rather than a CSS media query. This is a design decision to
make deliberately, not a port to grind through.

\item[\code{app/public/index.html} and \code{manifest.json}]
The entire PWA shell: \code{apple-mobile-web-app-*} metas, a script that rewrites the status-bar
style at launch, the \code{react-native-web} CSS reset, home-screen icons. All of it is replaced by
native configuration --- and the \code{black-translucent} status-bar hack and the iOS 26 top-blur
workarounds disappear with it. A simplification, not a loss.

\item[Focus and the keyboard, four files]
\code{features/auth/\bk focusField.ts} uses \code{document.getElementById().focus()}, which becomes a
\code{TextInput} ref. \code{features/ranking/\bk Compare.tsx}'s \code{window} \code{keydown}
shortcuts are simply inert on a phone. In \code{lib/useKeyboard\bk Inset.ts} and
\code{features/search/\bk keyboardLift.ts}, the \code{window.visualViewport} and
\code{innerHeight} compensation for iOS Safari's fixed viewport retires in favour of the native
branch (\code{Keyboard.\bk addListener('keyboardWillShow')}) already written beside it.

\end{findings}

\subsection{Everything else}

\begin{findings}

\item[\code{app/settings/privacy.tsx}]
\code{fetch('/privacy.md')} is a relative URL resolved against the Vercel origin. A native bundle has
no origin, so this fails outright: bundle the Markdown as an asset, or serve it from an absolute URL.

\item[\code{features/profile/\bk useShare\bk Profile.ts}]
\code{profileUrl()} returns an empty string off the web, and \code{Share.share} is then called with
it --- so the native branch, as written, shares nothing. It needs a canonical https base URL, plus
Associated Domains and an \code{apple-app-\bk site-association} file on the Vercel domain if shared
links should open the app rather than the website.

\item[\code{features/create-bev/\bk useOnline.ts}]
\code{navigator.onLine} and \code{window} online/offline events become
\code{@react-native-\bk community/\bk netinfo} --- a new dependency, and therefore its own small PR
under \code{app/CLAUDE.md}'s rules.

\item[\code{features/bug-report/\bk context.ts}]
\code{osFamily()} and \code{browserFamily()} parse a user-agent string, so both return
\code{undefined} on a device and \code{platform} degrades to the bare string \code{'ios'}. Send
\code{expo-device}'s OS version and the native build number instead --- a
\code{BugReport\bk Environment} contract change.

\item[\code{docs/find-bev-in-\bk app-plan.md} and the planned \code{MapView.web.tsx} (M-16)]
The Map tab and the Find Bev port are designed around the Google Maps \emph{JavaScript} API and
WebKit's sensor behaviour: one URL for the whole trip, a React Native \code{Modal} instead of a
route, and an iPhone spike (S2) gating the entire feature. Going native dissolves risks R3 and R6 ---
Home-Screen camera, motion and geolocation flakiness --- in favour of \code{react-native-maps},
\code{expo-camera}, \code{expo-sensors} and \code{expo-location}. \textbf{This is the strongest
engineering argument for the move}, and it is also why sequencing matters: it invalidates part of an
already-reviewed plan.

\item[\code{.github/workflows/app.yml} and \code{app/scripts/\bk check-bundle.mjs}]
CI typechecks, runs \code{jest-expo}, \code{expo-doctor} and the web export, and
\code{check-bundle.mjs} scans \code{dist/} to prove the jest-only \code{FakeApi} never ships.
\code{jest-expo} can run an \code{ios} project beside the web one; the leak guard needs a native
equivalent, since it only ever inspects the web \code{dist/}; and prebuild validation plus EAS
build and submit want a place in the workflow.

\end{findings}

% --------------------------------------------------------------------- 4 ---
\section{App Review and compliance}

\begin{warnbox}
\textbf{\color{accent}Blocker 1 --- user-generated content, guideline 1.2}

Bevi carries comments, photos, display names and profiles. Guideline 1.2 requires a reporting
mechanism \emph{and} a way to block abusive users. Our contract has \code{bevs.report(id, reason)}
and nothing else: TASKS.md records comment reports as deferred, and nothing anywhere in
\code{shared/src/api.ts} blocks a user. Expect to build block/mute plus comment reporting before
submission --- backend and app, both lanes. This is the item most likely to cost us a rejection
cycle, so it should be scheduled early rather than discovered at submission.
\end{warnbox}

\begin{warnbox}
\textbf{\color{accent}Blocker 2 --- the media contract}

\code{UploadMediaInput.file: Blob} cannot work from a device as written (\S\ref{sec:media}). It is a contract
change, so it touches \code{shared/}, \code{packages/api}, the \code{FakeApi} and the conformance
suite, and it crosses both lanes. Settle the shape before phase 1 rather than during it.
\end{warnbox}

\textbf{Alcohol.} Three alcoholic categories sit behind \code{canSeeAlcohol}. Under Apple's 2025
rating tiers (4+, 9+, 13+, 16+, 18+), alcohol references raise the rating; our existing server-side
gate is the right answer to give the questionnaire. It must keep holding on every new native surface
--- notably any Find Bev result that could contain a liquor store.

\textbf{Invite-gated sign-up (guideline 2.1).} App Store Connect's review notes need a working demo
account \emph{and} a valid invite code. Without both, the reviewer meets a wall and rejects.

\textbf{Already satisfied, and worth knowing.} In-app account deletion exists
(\code{settings/delete-account}), which meets guideline 5.1.1(v). And because auth is username and
password with no third-party social login, \textbf{Sign in with Apple is not required} --- a
genuinely non-trivial saving. App Privacy disclosures must still cover photos, the bug reporter's
evidence uploads (a purge script already exists) and, later, location.

% --------------------------------------------------------------------- 6 ---
\section{Three decisions to make first}

Each of these changes the work in the plan below, so they are worth settling first.

\begin{callout}
\textbf{\color{ink}1. Dual target, or native-only?}

Keeping the Vercel web build means every native-only API needs a \code{.web.tsx}/\code{.native.tsx}
split or a \code{Platform} branch, permanently. Dropping it retires \code{vercel.json},
\code{public/} and roughly a third of the inventory in \S\ref{sec:inventory} --- at the cost of shareable web links and
the browser as a quick demo surface.
\end{callout}

\begin{callout}
\textbf{\color{ink}2. Bev icons: bundled, or served?}

8.2\,MB inside the IPA, against a CDN dependency with a cold-launch cache to warm. Both are
defensible; the choice determines whether \code{scripts/icons/export-app.mjs} grows an asset-map mode
or an upload step.
\end{callout}

\begin{callout}
\textbf{\color{ink}3. Does Find Bev wait for native?}

Its reviewed plan is built for WebKit, and its iPhone spike S2 is the next step on the current path.
Going native first makes that spike largely moot and the feature materially easier to build. That
argues for ordering the platform move \emph{before} the Map tab, not after --- but it means
re-opening part of a plan that has already been through review.
\end{callout}

% --------------------------------------------------------------------- 5 ---
\section{Sequence and effort}

\vspace{0.4em}
\renewcommand{\arraystretch}{1.45}
\noindent\begin{tabular}{@{}p{0.085\textwidth}p{0.545\textwidth}p{0.25\textwidth}@{}}
\toprule
\textbf{Phase} & \textbf{Work} & \textbf{Rough effort} \\
\midrule
0 & Enrol; bundle id; the \code{ios} block in \code{app.config.ts}; \code{expo prebuild};
    \code{expo run:ios} on a real phone. \newline
    \textit{Done when: it launches and signs in on a device.} & 1--2 days \\
\addlinespace[0.35em]
1 & Session storage and \code{AppState} refresh; the media contract; the icon strategy; the privacy
    document; NetInfo; the photo pipeline. \newline
    \textit{Done when: every flow works on a device.} & 1--2 weeks \\
\addlinespace[0.35em]
2 & Accessibility sweep (the 83 ARIA props); haptics pass; real glass; sheet, status-bar and
    safe-area polish. \newline
    \textit{Done when: VoiceOver parity with the web build.} & 1--2 weeks \\
\addlinespace[0.35em]
3 & Block and comment reporting (both lanes); privacy manifest and disclosures; the rating
    questionnaire; the demo account. \newline
    \textit{Done when: accepted for external TestFlight.} & 1--2 weeks \\
\addlinespace[0.35em]
4 & CI: a native jest project, prebuild validation, EAS build and submit, a native leak guard.
    \newline \textit{Done when: \code{main} can cut a build.} & 2--3 days \\
\bottomrule
\end{tabular}
\renewcommand{\arraystretch}{1.0}

\vspace{0.7em}
{\color{mute}\small\noindent Measured on 2026-11-14 against \code{be/custom-bev-icons} @
\code{31381bc}; nothing in the repository was changed to produce this document.\par}

\end{document}
